Distributed software-defined networks (SDN) use several controllers, but time-varying traffic can overload some controllers while others sit idle and still consume energy. Static controller placement and purely reactive switch migration address only part of this problem. This project investigates whether a deep reinforcement learning agent can manage controllers proactively at runtime. A Rainbow Deep Q-Network (RDQN), acting as a master controller, selects at each decision interval one of four action types: migrate a switch away from an overloaded controller, evoke a parked controller, park an underloaded controller, or take no action. A valid-action mask with feasibility checks keeps every chosen action executable, and a bounded reward combines switch--controller latency, modelled controller energy and traffic load balance. An optional gated recurrent unit (GRU) forecaster adds predicted loads to the state so that capacity can be added before overload occurs.

The framework was evaluated on five Internet Topology Zoo networks (9--197 switches) with chronologically split synthetic traffic, and on Abilene-derived traffic mapped to OS3E. Each learned policy was compared with threshold, random-valid and no-action baselines under identical conditions, and each test split was evaluated once after model selection. On held-out synthetic traffic, the selected learned policy achieved the highest reward on all five topologies and used 4.2--19.6\% less simulated energy than the threshold policy, with slightly higher average latency on four topologies. Forecasting helped on OS3E, improving reward from $-478.46$ to $-452.94$ and lowering energy from 442.18\,W to 380.89\,W relative to current-state RDQN, but did not help on Gridnet or BellCanada. In the real-traffic study, predictive GRU-RDQN had the better mean reward over three seeds ($-111.19\pm17.53$ versus $-121.80\pm20.14$), yet current-state RDQN won in two of the three seeds. Finally, the unchanged OS3E policies were run on a Mininet emulation with Open vSwitch and five separate os-ken OpenFlow controllers, replaying the same traces over three 150-step windows. All 687 controller start, stop and switch-handover operations succeeded, and GRU-RDQN made exactly the same decisions as in the simulator in every window, achieved the best reward in all three predictive-mode runs, and was unaffected by small measurement noise; a switch handover took 40\,ms and a controller start 0.65\,s on average. The contributions are a runtime migrate/park/evoke formulation, a leakage-controlled evaluation protocol, evidence that the value of forecasting is scenario-dependent, and an emulation backend that executes the learned decisions on a real OpenFlow data plane.
